% Generated by D:/tools/scratch_qdd/mmrec/render.py from data_e.py - edit the data, not this file.
\subsection{E. 데이터 --- 자체 시뮬이 바닥, 공개 데이터는 픽셀로 바꿔 쓰고, 다양성은 필수}\label{rec:E}
\noindent\textbf{한 줄:} 조종 표본·복구·자기 인식은 공개 대체가 없어 자체 시뮬(참값 모델)이 주원천이다. 남의 로봇 데이터는 행동 라벨이 아니라 픽셀·QA로 바꿔 쓰는 것이 문헌의 방식이고, 지금 과제는 같은 탁자·같은 물체라 다양한 데이터는 필수(OOD 평가 필수). 공개 데이터 변환은 시제품·관문까지, 좌표계 없는 원천은 픽셀 경로(T0 + T5)가 주 추천. 8B 검증 등록(E-OPEN8·E-XEMB8)은 초안.

\paragraph{E1. 자체 시뮬 R2와 소규모 실물 S-E2E}
{\small 09-24 \textasciitilde{} 09-26 09:53 KST}
\begin{itemize}
  \item \textbf{사용자} (ul 66, 09-25 04:18 KST): \intent{직므은 단순이니까 10으로 하되 나중에는 우리가 데이터하고 할떄는 30으로 모을거}
  \item \textbf{그래서(결정):} §62: S-E2E(ROBOTIS 공개 RB1·RB2)는 10 Hz 그대로, 우리 수집 데이터는 30 Hz. §66: R2 생성기 확정(30 Hz·다과제·도메인 무작위화·LeRobot).
  \item \textbf{결과:} R2\_TRAIN(b2b0b3e 08:30, LeRobot 18d4e2f 09:53): 6,000편 생성, 유효 5,126(85.4 \%) --- mug\_tray 99.1 \%·mug\_marker 98.7 \%·bottle\_tray 58.5 \%(실패 874 중 794가 병 놓기); 단계 B 행 1,495,348, LeRobot v2.1 6개 22.13 GB, verify 6/6. S-E2E 42,813행(D1).
  \item \textbf{참고:} 이 데이터가 E-SR1c(시뮬 분기)·E-TEACH-L8(TRAIN 시드)·E-PT의 바탕이다. 지금 시뮬 과제는 탁자 높이 z 0.850·같은 탁자·같은 물체·방해물이다(E6의 출발점).
\end{itemize}

\paragraph{E2. ROBOTIS AI Worker 공개 데이터 전수 조사}
{\small 09-26 15:29 -- 18:43 KST}
\begin{itemize}
  \item \textbf{사용자} (ul 105, 15:29 KST): \intent{1. ai워커에 들어가보면 로보티즈 오픈 데이터 많은데 그중에 쓸만한거 찾아봐 2. 고}
  \item \textbf{해 봄:} 조사(1ef9f2c, 16:20, 학습·유료 0): HF의 FFW/AI Worker 관련 1,150개 데이터셋 메타데이터 전부(ROBOTIS 조직 7개 + 직원 계정 RobotisSW·Dongkkka + 외부), 상위 14개는 프레임 눈 확인. (`2. 고' = E-MAR-real 진행, D4.)
  \item \textbf{결과:} 1순위 \texttt{Dongkkka/\allowbreak{}ffw\_\allowbreak{}bg2\_\allowbreak{}rev4\_\allowbreak{}pickup\_\allowbreak{}obj\_\allowbreak{}1127\_\allowbreak{}total2}(806편, 우리 기종 BG2, 카드 apache-2.0; 같은 공구 벽에서 지시문이 대상을 고르는 쌍 601·같은 물체를 다른 팔로 457), 2순위 RobotisSW 과제 카탈로그(322과제·7,364편, apache-2.0, 지시문은 번호뿐), 3순위 Merged\_PickPlace 병/캔(두 대상, 라이선스 없음). 같은 팔·같은 장면의 대체 경로 시연은 어디에도 없음; 실패 데이터는 외부 사용자 것뿐.
  \item \textbf{결과:} 2차(66f2834, 18:43, 사용자 후속 질문 `VLA 오픈 데이터'): 공식 AI Worker VLA 데이터 = HF ROBOTIS 7개(+ GR00T 미세조정 모델 6개). 약 10,500편 FK 채굴: 같은 시작·같은 목표에서 $\geq$ 5 cm·$\geq$ 40 px 갈라지는 쌍 RB2 420·bg2 192·pickup\_obj 약 900; $\geq$ 10 cm·$\geq$ 80 px는 라이선스 있는 원천에서 약 100쌍뿐. 프레임상 의도적 우회가 아니라 사람의 좌우 흔들림.
  \item \textbf{그래서(결정):} 이 쌍들은 오프라인 궤적 추종 평가(n $\geq$ 300)·두 대상 평가(pickup\_obj 601 + RB2 426)에만 쓰고 MolmoAct 조종 학습 데이터로는 쓰지 않는다. 1순위 데이터는 RB3로 추가(E3).
\end{itemize}

\paragraph{E3. RB3 추가 --- 새 데이터 판 se2e\_c2}
{\small 09-26 18:52 -- 18:56 KST}
\begin{itemize}
  \item \textbf{결과:} 2f973ba(18:56): RB3 806편·104,470프레임·3.31 GB(verify 오류 0) $\rightarrow$ 변환 627편·16,721행(그리퍼를 안 쓴 빈 편 179 제외, 프레임 확인), train 15,906 / val 815. RB1·RB2와 해시 일치 0·궤적 최근접 0.351 rad. 조사와 다른 사실: 손목 영상 회전 불필요, 파켓 task\_index 58편 오류(episodes.jsonl 사용), \texttt{release\_\allowbreak{}visible} 240/627(눈 대조 24/24).
  \item \textbf{그래서(결정):} \texttt{se2e\_\allowbreak{}c2} = \texttt{se2e\_\allowbreak{}c1} + RB3(RB1·RB2 행은 바이트 동일). 카드는 apache-2.0이지만 직원 개인 계정이라 내부용 표시. E-SR1e는 일반화 차가 없어 RB3를 팔에 넣지 않았다(D9).
\end{itemize}

\paragraph{E4. 상위 VLM 학습 데이터 조사 --- 시뮬이면 더 좋다}
{\small 09-26 21:0x -- 21:17 KST}
\begin{itemize}
  \item \textbf{사용자} (ul 111, 21:0x KST): \intent{저거에 맞는 데이터를 찾는 작업을 꼼꼼하게 해보자 심이면 더 좋을 것 같구}
  \item \textbf{결과:} 조사(73787c6, 21:17): 약 40개 데이터셋 확인, 7종 파드 표본 확인. 판정: 조종 표본·복구(공개 수치 교정 데이터 없음)·자기 인식(AI Worker 공개 시뮬 데이터 0건)은 자체 Isaac 시뮬 + 참값 모델이 주, 공개는 장면 다양성·일반 그라운딩 보조. 상위 5: MolmoBot-data(RBY1)·BEHAVIOR-1K 2025·RoboTwin 2.0·RefSpatial·InternData-M1(게이트·NC). 막힘: 게이트 동의, NC·ND 라이선스는 Astra 전송 보류.
\end{itemize}

\paragraph{E5. 남들은 맞지 않는 로봇 데이터를 어떻게 쓰나 --- 35B 데이터 양}
{\small 09-26 22:2x -- 22:29 KST}
\begin{itemize}
  \item \textbf{사용자} (ul 115, 22:2x KST): \intent{기존 데이터를 쓸 방법을 찾아봐 다른 연구들에서는 우리와 같지는 않겠지만, 다른사람들은 자기에 맞지 않은ㄴ 로봇데이터를 어떤식으로 튜닝 혹은 활용을 했는지 찾아봐 35b학습시키는데 데이터가 충분해야하다고 생각하거든 거기다 범용성도 있으려면 말이야}
  \item \textbf{결과:} 조사(2373a02, 22:29): 남들은 맞지 않는 로봇 데이터를 행동 라벨로 쓰지 않고 (a) 기체 무관 영상 표현(MolmoAct 그리퍼 2D 궤적 + 깊이 토큰, Embodied-R1 점, HAMSTER[기초]) 또는 (b) QA·추론 재라벨(Robo2VLM, InternVLA-M1, RoboBrain 2.0)로 바꾸고, 출처 표시(X-VLA)·일반 데이터 5--25 \% 공동학습으로 범용성을 지킨다. Vlaser(ICLR 2026): 남의 영역 데이터만으로는 조종 이득 없음(41.8 $\rightarrow$ 43.2 \%, 같은 영역 65.1 \%).
  \item \textbf{그래서(결정):} 자체 시뮬 주 판정 유지, 공개는 픽셀·카메라 좌표 인식 QA로(출처·좌표계·내부값을 글로). 35B 1차 약 40만 표본(자체 조종 25 \%·자체 인식 25 \%·공개 인식 25 \%·추론 10 \%·일반 공간 10 \%·일반 5 \%) [추정]. 무료 검증 E-XEMB8은 설계만, L8 결과 뒤.
\end{itemize}

\paragraph{E6. 다양한 데이터는 필수 + OOD 평가}
{\small 09-26 23:1x -- 23:36 KST}
\begin{itemize}
  \item \textbf{사용자} (ul 116, 23:1x--23:2x KST): \intent{근데 쓸만한 거는 가공해서 조종명령 답하게 바꾸는건 더 좋은건 아니야?}
  \item \textbf{사용자} (ul 116, 23:1x--23:2x KST): \intent{왜 그러냐면 우리의 테스크는 동인 높이 책상에 동일한 장애물들에 다 동일하잖아 무조건 다른데이터가 필수로 있기는 해야해}
  \item \textbf{그래서(결정):} 정본 §97 보충 1(94c0440, 23:36): 상위 VLM 학습에서 다양한 데이터는 필수(넣을지를 실험으로 정하지 않음, E-XEMB8은 형태·비율만). (a) 공개 변환 격상 --- 인식 QA + 좌표계 명시 조종 C$'$(\texttt{source}·\texttt{frame: base\_\allowbreak{}<기체>}, 그 기체 자기 정보 블록 자동 생성, 그리퍼 사건에서 우리 단계 어휘로 자름); (b) 자체 생성기 다양화 --- 탁자 높이(목표 $\pm$10--15 cm)·크기·색, MolmoSpaces 물체, 방해물 0--6·혼동체, 받침, 조명·질감, 카메라 장착 흔들림(머리 각도는 고정), 새 과제 변형 4; 기본 경로는 바이트 동일.
  \item \textbf{그래서(결정):} OOD 평가(학습 절대 금지): 보호 분할 \texttt{ood} 시드 70000--70999 --- OOD-H 높이·OOD-O 물체 범주 20 \% 보류·OOD-D test 풀 방해물·OOD-X 공개 기체 보류; 지표는 L8과 같고 프레임 누수율 0이 관문.
  \item \textbf{진행 중·대기:} \tentative{사용자 확인 필요로 남은 것: 탁자 높이 범위를 리프트(\texttt{set\_\allowbreak{}lift})로 맞춰도 되나, 게이트 데이터 HF 동의 여부. 생성기 다양화·OOD 세트는 [예정].}
\end{itemize}

\paragraph{E7. 공개 데이터 변환 방법 --- MolmoAct 방식만으로 되나}
{\small 09-27 00:14 KST (커밋 시각)}
\begin{itemize}
  \item \textbf{사용자} (ul 119, 00:3x KST): \intent{1번을 할 방법을 더 자세하기 찾아봐야해}
  \item \textbf{사용자} (ul 119, 00:3x KST): \intent{몰모엑트가 한방식데로 데이터를 하면 안되나?}
  \item \textbf{사용자} (ul 119, 00:3x KST): \intent{우리는 몰모엑트방식으로는 데이터가 부족한가?}
  \item \textbf{결과:} \texttt{public\_\allowbreak{}data\_\allowbreak{}conversion\_\allowbreak{}howto} + CPU 시제품 \texttt{tools/\allowbreak{}xemb/\allowbreak{}}(624fdd6, 00:14; 시험 31개, 유료 0·GPU 0): MolmoAct 공개분(CC BY, 약 1,000만 프레임)은 픽셀 점·궤적은 넘치지만 기저 좌표 조종(C$'$)·metric 3D·평면은 0이라 그것만으로는 부족; 보정 있는 원천은 포인팅 대신 FK + 보정 투영. C$'$ 주원천 MolmoBot Franka(위에서 잡기 93 \%) $\rightarrow$ RoboTwin(73 \%) $\rightarrow$ BEHAVIOR(29 \%). 시제품 관문: EE가 그리퍼 위 BEHAVIOR 37/37·MolmoBot 자동 90.2 \%(눈 $\geq$ 97 \%)·RoboTwin 눈 9/9, 단계 $\leftrightarrow$ 계획기 94.2 \%, 탁자 평면 61/61, 분할 $\cap$ 깊이 3D 중심 9.1 cm 불통.
  \item \textbf{그래서(결정):} 3D 답은 참값(GT)에서만 만든다. 변환 결과 P 1,145·C$'$ 54(시제품). 세트 구성은 E-OPEN8·E-XEMB8 초안으로(E9).
  \item \textbf{참고:} user-log 119의 시각 표기는 15:3x UTC(= 09-27 00:3x KST)이고 커밋은 00:14 KST라 순서가 맞지 않는다 --- 기록 그대로 둔다.
\end{itemize}

\paragraph{E8. 좌표계 없는 공개 데이터 --- 가장 효율적인 방법 채택}
{\small 09-27 00:4x -- 00:49 KST}
\begin{itemize}
  \item \textbf{사용자} (ul 124, 00:4x KST): \intent{제일 효율적인 방법으로 채택을 해야해 그것도 연구 자료 뒤져서 저 세가지 방법외에 하나 더 찾아봐 2개도 괜찮고}
  \item \textbf{사용자} (ul 124, 00:4x KST): \intent{너가 아까 말한 좌표계에서 픽셀로만 내는것도 방법이 될 수는있고}
  \item \textbf{사용자} (ul 124, 00:4x KST): \intent{근데 저 4개의 방법론은 오픈소스데이터로 일단 해야해 이해했으?}
  \item \textbf{사실 확인:} 기존 트랙(다른 에이전트, \texttt{tools/\allowbreak{}xemb/\allowbreak{}routes.\allowbreak{}py}): T0 픽셀만, T1 좌표계 만들기(FK 3D + 검출 2D $\rightarrow$ PnP), T2 암묵 학습, T3 단안 metric 깊이.
  \item \textbf{결과:} 조사(a455fc8, 00:49, 코드·GPU·유료 0): 새 방법 T4 로봇 몸 정합 보정(CtRNet-X ICRA 2025 --- DROID 공식 2025 재보정 약 3.6만 장면에 쓰임, 렌더 IoU 0.019 $\rightarrow$ 0.836), T5 점 추적 픽셀 감독(CoTracker3 ICCV 2025, Magma CVPR 2025가 OXE 약 97만 궤적 라벨에 사용) + 관절값 그리퍼 사건으로 집기·놓기 점. 스테레오·다시점 재구성은 T3 변형, 잠재 행동·Cloak은 뺌.
  \item \textbf{그래서(결정):} 정본 §97 보충 3(추천안, 검증 후 확정): 보정 없는 공개 데이터(AI Worker·OXE)에 먼저; 주 = 픽셀 경로 T0 + T5(E-PT \texttt{point} 형식 C-pt 포함, m 값 없음), 보조 = T1(+ T4, 재투영 $\leq$ 5 px 관문 통과분만 투영·C$'$), T3는 관문 전 서열 QA만, T2는 방법이 아니라 결과(ND-1로 판정). 공유 앞단은 추적(T5).
  \item \textbf{진행 중·대기:} \tentative{확정 절차: L단계(라벨 정확도 --- 보정 있는 공개 데이터는 보정을 가린 참값, AI Worker·OXE는 클릭 약 1,000점) $\rightarrow$ T단계 8B 팔 A·PX·CAL(·D3), E-PT와 같은 DEV·OOD-H + OOD-X-open. 둘 다 아직 실행 전.}
\end{itemize}

\paragraph{E9. 8B 검증 등록 초안 E-OPEN8·E-XEMB8}
{\small 09-27 KST (미커밋 초안)}
\begin{itemize}
  \item \textbf{진행 중·대기:} \tentative{\texttt{docs/\allowbreak{}stage3/\allowbreak{}prereg\_\allowbreak{}open8.\allowbreak{}md}·\texttt{prereg\_\allowbreak{}xemb8.\allowbreak{}md}는 작업 트리에만 있는 초안(`등록 아님, 실행 안 함')이고 커밋되지 않았다.}
  \item \textbf{참고:} E-OPEN8 초안: O-A 공개 데이터만으로 LoRA한 8B가 우리 요청 형식에서 대상에 다가가나(영샷 대비), O-B L8 자체 데이터 + 공개 세트를 같은 걸음 수로 섞으면 DEV·OOD-H가 나아지나; 좌표 누수 0; E-PT와 같은 모델·816걸음·같은 스냅숏·같은 채점기; 팔당 약 1.4 GPU-h [추정].
  \item \textbf{참고:} E-XEMB8 초안(§97 보충 1 개정판): 같은 상태·같은 계산량에서 L8 학습 행 일부를 (B) 공개 인식 QA, (D) 공개 C$'$, (E) 둘 다로 바꿔 가르기; 선행 = E-PT 결과(인터페이스)·관문 통과 세트. 3팔 약 4.2 GPU-h [추정].
\end{itemize}

\paragraph{E10. 좌표계 없는 오픈소스 --- 세 트랙 구현, 오픈소스 먼저, 점 추적 공용}
{\small 09-27 01:0x -- 02:5x KST}
\begin{itemize}
  \item \textbf{사용자} (ul 123, 01:0x -- 02:5x KST): \intent{아니 아까부터 계속 이얘기 했잖아 좌표계 없는거는 어떻게 할건지 오픈소스를 어떻게 할건지 해야한다니까?}
  \item \textbf{사용자} (ul 123, 01:0x -- 02:5x KST): \intent{좌표계가 없는걸 좌표계를 만들건지 / 그냥 이미지를 엄청 ㅁ낳이 핛브햇거 좌표계를 뭔가 너머에서 알 수 있게 하던지 / 뎁스를 근사화하는 비슷한 방식으로 해서 하던지}
  \item \textbf{사용자} (ul 123, 01:0x -- 02:5x KST): \intent{일단 오픈소스를 가지고 학습을 할건데 어떤걸 써서 어떻게 할거임? 대규모는 아니고 소규모로 해두돼}
  \item \textbf{사용자} (ul 123, 01:0x -- 02:5x KST): \intent{점추적은 T,1,3에 다 들어갈 수 있을 것 같은데 아닌가?}
  \item \textbf{결과:} \texttt{tools/\allowbreak{}xemb/\allowbreak{}} + howto 9--12절(7faca71): RB2·MolmoAct OXE에 먼저, 보정 있는 공개 세트는 보정 숨긴 정답 대조만. 공용 점 추적 앞단(CoTracker3). T1: 정확한 2D면 0.81 px, 포인팅·추적 2D면 24--26 px(불통). T3(MoGe-2): 물체 17.5 cm·평면 7.7 cm(불통). E-OPEN8 오픈만 6,500행.
\end{itemize}

\paragraph{E11. 작전 T --- 깊이 추정은 버리고 두 안(T1+T4 대 C$'$)을 끝까지 키운 뒤 결정}
{\small 09-27 03:0x KST \textasciitilde{} (최종 결정 12:53 KST)}
\begin{itemize}
  \item \textbf{사용자} (ul 144, 03:0x KST): \intent{일단 내가 시간을 더 줄테니까 계속 도전해보다가 10시간 줄게 안되면 잘되는 탑2에 다 몰아주고 그다음에 탑1에 다 몰아주는걸로 할게}
  \item \textbf{사용자} (ul 144, 03:0x KST): \intent{아니다 깊이 추정은 그냥 아닌게 맞는 것 같아 버리고 2개중에 고민하고 계속 각각을 최대로 발전시킨 후에 10시간 뒤에 최종 결정하자. 이걸 작전명 T라고 할게}
  \item \textbf{결과:} T4 실루엣 정렬(d2b2251·caea370): RB2 보류 경계 2.32 px, 프레임 96 \% $\leq$ 5 px $\rightarrow$ RB2 투영 점 18,636. 명목 카메라도 4.65 px라 T1의 21--42 px는 포인팅 잡음. RB3 2.87 px(73 \%), RB1 5.48 px(44 \%, 팔 마스크 병목)(5033591). C$'$ 물체 이름을 과제 문장에서(`the object' 누수 제거, 746396c). T3 재검토: MoGe-2 + FK 척도 정렬도 물체 19 cm·평면 12 cm $\rightarrow$ 폐기 유지. E-OPEN8 O-A 옛 시뮬 예비: DEV 106.8 mm(영샷 0.85배, NONE), 누수 1/305.
  \item \textbf{진행 중·대기:} \tentative{최종 결정 09-27 12:53 KST. 선택은 L8-X에서 같은 조건 짝 비교 뒤(ul 143). 논문 본문 3.3절·5절에 `결정 대기'로 반영.}
\end{itemize}

\paragraph{E12. E-DIST8 --- 데이터 3수준 $\times$ 공개 픽셀 유무 6팔}
{\small 09-27 03:3x KST}
\begin{itemize}
  \item \textbf{사용자} (ul 141, 03:3x KST): \intent{A·B·C·D 모두에 픽셀을 넣으면 픽셀이 도움인지 방해인지 확인하는 것까지 넣는걸로 하자}
  \item \textbf{그래서(결정):} E-DIST8(8B, 무료, 같은 걸음·표본 수): A(L8)·A+px, B(L8-D)·B+px, C(L8-D + 공개 3D: BEHAVIOR 계량 3D QA + C$'$)·C+px. 평가 DEV·OOD-H(좁게·넓게)·OOD-O/D. 픽셀 효과는 세 수준 각각 짝 비교. L8-D 준비 뒤 등록. 논문 6절 `이번 마일스톤'에 반영.
\end{itemize}

\paragraph{E13. L8-X 환경 모음 --- 시뮬을 약 10배 다양하게, 데이터셋은 1.5년 규칙 해제}
{\small 09-27 03:5x -- 04:5x KST}
\begin{itemize}
  \item \textbf{결과:} L8-D 생성기 등록(e4dced8): 한 Isaac 프로세스에 한 높이, 높이별 작업 상자 = 머리 시야 띠 $\times$ IK 도달 띠, 관문 G-H --- 고정 리프트 0.74--0.98 m 통과. L8-X 설계(dfeb965): 받침면 6(탁자·조리대·계단식 선반·낮은 탁자·바닥 통/상자·윗판 수납장 = 한계 탐침), 과제 약 12, 방해물 0--8, 외관 확장. 물체·과제(e6b2e1e): 받침대·파란 머그·작은 빨간 컵·회색 통, 학습 X 과제 7종. 보류 분할 학습 전 고정(d37634b): OOD-O 작은 컵, OOD-T 보류 조합 2종, OOD-S는 가구 장면 들어올 때. 여러 단계 과제 2종(aeb37f3).
  \item \textbf{사용자} (ul 147, 04:5x KST): \intent{데이터는 1년반해제해줄게}
  \item \textbf{그래서(결정):} CLAUDE.md 예외 줄(99a98d9·d0fdc10): 데이터셋에는 1.5년 규칙을 적용하지 않음(DROID·RH20T·OXE 등 사용 가능), 논문·방법 근거에는 유지; 라이선스·게이트·신뢰도·Astra 전송 제한은 그대로. 논문 3.3절에 한 문장.
\end{itemize}
